\section{Projective completion and literal restriction}
This section counts what survives when the standard projective star
circuit is restricted to its affine subgraph, and compares that count
with a star circuit built directly for the affine graph.
\begin{proposition}[Literal restriction of the line-star circuit]
Let $q$ be a prime power and $N=q^2+q+1$.  The line-centred star
circuit for the incidence edge function of $\operatorname{PG}(2,q)$
has $N(q+2)-1$ gates.  Set the variable of one line $\ell_\infty$
and all its $q+1$ point variables to zero.  Propagating zeros in this
circuit leaves exactly
\[
 q(q+1)^2-1
\]
gates, which is $q$ more than the affine point-star upper bound
$q^2(q+2)-1$.  The literal deletion count is $q^2+2q+2$.
\end{proposition}
\begin{proof}
Each projective line has a neighborhood OR tree with $q$ gates,
followed by one AND; the outer OR tree has $N-1$ gates.  Every line
other than $\ell_\infty$ loses exactly its unique infinity point.
Its neighborhood tree therefore retains $q-1$ OR gates and its AND.
There are $N-1=q(q+1)$ surviving lines and $N-2$ outer OR gates, giving
\[
 (N-1)(q-1)+(N-1)+(N-2)=q(q+1)^2-1.
\]
Subtracting from $N(q+2)-1$ gives $q^2+2q+2$.
The affine point-star circuit has $q^2$ stars, each with $q$ neighborhood
ORs and one AND, and $q^2-1$ outer ORs.  Its size is $q^2(q+2)-1$,
smaller than the inherited circuit by $q$.
\end{proof}

